In many collective decision-making tasks, a group of agents (\eg, voters) needs to select a subset of alternatives (\eg, candidates) that provides equitable outcomes for each agent while respecting additional demographic fairness constraints.
We consider this problem in the context of committee selection, which generalizes the single-winner selection problem to selecting multiple representatives.
Specifically, we assume that agents and alternatives are embedded in a metric space, where the distance between an agent and an alternative represents the cost (or dissatisfaction) that the agent experiences when that alternative is selected. 
Each agent can provide a complete ranking over the alternatives that is consistent with these underlying distances: if an agent ranks alternative $a$ above alternative $b$, then the distance from the agent to $a$ is no greater than the distance from the agent to $b$. 
The computational social choice literature refers to these preference rankings as ordinal information, as opposed to the explicit distance values, which are referred to as cardinal information. 
The goal is to select a committee (subset of alternatives) that minimizes the social cost, defined as a function of the distances of agents to their closest member in the selected committee. 

While ordinal rankings are readily available, the exact distance values are typically unknown to the algorithm.
A common way to evaluate the performance of a social choice algorithm that uses only ordinal information is through the notion of \emph{distortion}, first introduced by \citet{procaccia2006distortion}.
This concept was later extended to metric preferences by~\cite{anshelevich2018approximating}, which measures the quality of the solution produced by the algorithm, defined as the worst-case ratio between the social cost of the algorithm's solution and the social cost of the optimal solution that has full access to all pairwise distances.
In recent years, metric distortion has been a central tool for evaluating the performance of social choice algorithms in the single-winner setting~\cite{anshelevich2018approximating, munagala2019improved, gkatzelis2020resolving, kizilkaya2022plurality, kizilkaya2023generalized, 
charikar2022metric, fain2019random,
charikar2024breaking}, \ie, when the committee consists of a single member.


Recently, \citet{pulyassary2025constant} and \citet{burkhardt2024low} extend the metric distortion framework from single-winner selection to $k$-committee selection, for arbitrary $k$. 
Their works consider both the utilitarian objective, which minimizes the sum of distances of each agent to its nearest committee member (which corresponds to the $k$-median or $k$-means objective), and the egalitarian objective, which minimizes the maximum distance of an agent to its nearest committee member (which corresponds to the $k$-center objective).
\citet{burkhardt2024low} show that for $k \geq 3$, no algorithm can achieve a constant distortion for the ordinal $k$-committee selection problem with the egalitarian objective (i..e, when cardinal information is not available).
In light of this hardness result, both works additionally allow limited access to cardinal information through distance queries in order to obtain constant-factor distortion guarantees.



In this paper, we continue this line of work by studying a fair version of $k$-committee selection. 
Specifically, we study the problem of \emph{ordinal fair $k$-committee selection} with an emphasis on both \emph{equity} and \emph{fairness}. 
The equity principle ensures that no agent is unduly disadvantaged by the collective outcome; we encode this through the egalitarian objective, which minimizes the maximum distance from any agent to its nearest committee member.
To enforce fairness in the committee, we assume that the alternatives are partitioned according to an underlying demographic profile. The goal is to guarantee minimum representation from each demographic group in the selected committee. Numerous studies in social welfare highlight that such minimum representation can promote equity, inclusion, and informed policy decisions. For instance, quotas have been widely considered as mechanisms to secure diverse representation and balance opportunities across groups~\cite{barnes2020gender, croissant2019ethnic}. Our fairness constraints are designed to align with this notion of demographic fairness. Additionally, similar to~\cite{burkhardt2024low,cembrano2025metric}, we focus on the setting where alternatives and agents coincide.


Informally, the ordinal fair $k$-committee selection problem is defined as follows. 
We are given a set of agents embedded in an unknown metric space, partitioned into demographic groups. 
Each agent provides a complete ranking over all other agents that is consistent with the underlying distances in the metric space.
The goal is to select a committee of size $k$ using only a limited number of distance queries, subject to fairness constraints that ensure minimum representation from each group, while minimizing the maximum distance from any agent to its nearest committee member (the egalitarian objective).
The key difference between our work and \citet{pulyassary2025constant} and \citet{burkhardt2024low} is that we incorporate fairness constraints on minimum representation requirements for each demographic group.

When full cardinal information is available, the ordinal fair $k$-committee selection problem is equivalent to the fair $k$-center problem~\cite{kleindessner2019fair}. 
To develop algorithmic solutions for the ordinal fair $k$-committee selection problem, we build on techniques from the fair $k$-center problem. 
For consistency with existing literature, we formulate our problem as the ordinal fair $k$-center problem with limited cardinal information, providing a natural extension to clustering methodologies that maintains alignment with prior work. 
We give a formal definition of the problem in Definition~\ref{def:ordinal-fair-k-center-with-limited-distance-queries}.
